\begin{remark}
\label{remark-cohen-bound-cardinality}
Let $R$ be a ring. Let $\kappa$ be an infinite cardinal.
By applying
Example \ref{example-oka-family-bound-cardinality} and
Proposition \ref{proposition-oka}
we see that any ideal maximal with respect to the property of not being
generated by $\kappa$ elements is prime. This result is not so
useful because there exists a ring for which every prime ideal
of $R$ can be generated by $\aleph_0$ elements, but some
ideal cannot. Namely, let $k$ be a field, let $T$ be a set whose
cardinality is greater than $\aleph_0$ and let
$$
R = k[\{x_n\}_{n \geq 1}, \{z_{t, n}\}_{t \in T, n \geq 0}]/
\bigl(x_n^2\ (n\geq1),\ z_{t,n}^2\ (t\in T,n\geq0),\
x_nz_{t,n}-z_{t,n-1}\ (t\in T,n\geq1)\bigr)
$$
This is a local ring with unique prime ideal
$\mathfrak m=(x_n\mid n\geq1)$. Indeed, every prime contains all
$x_n$ and $z_{t,n}$ because their squares vanish, and the quotient
by all these elements is the field $k$. The relations
$z_{t,n}=x_{n+1}z_{t,n+1}$ put every $z_{t,n}$ in $(x_n)$.
Every element of this ideal is nilpotent: it uses only finitely many
of the square-zero generators, and a product of more factors than
the number of those generators has a repeated factor and vanishes.
The augmentation to $k$ shows the ring is nonzero.
Thus $\mathfrak m$ is the unique prime and is countably generated.

\medskip\noindent
The ideal $J=(z_{t,n}\mid t\in T,n\geq0)$ in fact requires exactly
$|T|$ generators. To prove the lower bound we construct maps detecting
each original $T$-label. For $N\geq0$ let
$$
C_N=k[X_1,\ldots,X_N,W_N]/(X_1^2,\ldots,X_N^2,W_N^2).
$$
Define $C_N\to C_{N+1}$ by $X_i\mapsto X_i$ for $i\leq N$
and $W_N\mapsto X_{N+1}W_{N+1}$.
The square-free monomials
$X_E=\prod_{i\in E}X_i$ and $X_EW_N$, for $E\subset\{1,\ldots,N\}$,
form a $k$-basis of $C_N$.
Their images are respectively $X_E$ and $X_EX_{N+1}W_{N+1}$,
distinct members of the corresponding basis of $C_{N+1}$.
Thus all these maps are injective. Put $C=\colim_N C_N$.
Concretely, take classes of pairs $(N,c)$ with $c\in C_N$, identifying
two pairs when their images agree at a common later stage. Addition
and multiplication are computed at a common stage; compatibility of
the ring maps makes them independent of the stage and representatives.
The injectivity of every transition map makes each $C_N\to C$ injective.
In $C$ the original classes satisfy
$W_{n-1}=X_nW_n$ for $n\geq1$, and
$W_0=X_1\cdots X_NW_N\ne0$ at every stage and in the colimit.

\medskip\noindent
For each $t_0\in T$, the assignments
$$
\rho_{t_0}:R\longrightarrow C,\qquad
x_n\longmapsto X_n,\qquad
z_{t,n}\longmapsto
\begin{cases}W_n,&t=t_0,\\0,&t\ne t_0\end{cases}
$$
respect every displayed defining relation, including its index range,
and hence define a ring map. In particular
$\rho_{t_0}(z_{t_0,0})=W_0\ne0$.
Suppose a family $(j_\lambda)_{\lambda\in L}$ with $|L|<|T|$
generated $J$. Each $j_\lambda$ is a finite sum of multiples
of generators $z_{t,n}$, so the labels occurring in all these
finite expressions form a set $T'\subset T$ of cardinality at most
$\max(|L|,\aleph_0)<|T|$; here $T$ is uncountable.
Choose $t_0\in T\setminus T'$. Then $\rho_{t_0}$ kills every
$j_\lambda$, irrespective of its coefficients, but does not kill
$z_{t_0,0}\in J$, a contradiction.
This proves the lower bound. The displayed generating set for $J$
has cardinality $|T\times\mathbf Z_{\geq0}|=|T|$, proving equality.
In particular $J$ is not countably generated.
More generally, for any prescribed infinite cardinal $\kappa$, choosing
$|T|>\kappa$ gives a ring whose only prime ideal is countably
generated, hence generated by at most $\kappa$ elements, while $J$
cannot be generated by $\kappa$ elements.
\end{remark}